% !TEX root = /home/benedict/Documents/Studium/Studienarbeit/Studienarbeit_Graphen_Metriken_Aehnlichkeit/Arbeit/T3101_Studienarbeit_GraphenMetriken_Weichselbaum_TINF2018.tex

\documentclass[a4paper,12pt,ngerman,chapterprefix=false,listof=totoc,bibliography=totoc]{scrreprt}

\usepackage[a4paper,left=2.5cm, right=2.5cm, top=2cm, bottom=2.5cm]{geometry}
\usepackage[ngerman, english]{babel}
\usepackage{blindtext}
\usepackage{helvet}
\usepackage{subcaption}
\usepackage[utf8]{inputenc}
\renewcommand{\familydefault}{\sfdefault}

\usepackage{microtype}
\usepackage{graphicx}
\usepackage{wrapfig}
\usepackage{enumitem}
\usepackage{amsmath}
\usepackage{index}
\usepackage{csquotes}
\usepackage[onehalfspacing]{setspace}
\usepackage{listings}
\usepackage{changepage}
\usepackage{acronym}
\usepackage{listings}
\usepackage{abstract}
\usepackage{scrhack}
\usepackage{booktabs}
\usepackage{pdflscape}
\usepackage{hyperref}
\usepackage{svg}
\hypersetup{
    colorlinks,
    citecolor=black,
    filecolor=black,
    linkcolor=black,
    urlcolor=black
}
\usepackage[bottom]{footmisc}

% Biblatex includes
\usepackage[style=alphabetic,backend=biber]{biblatex}
\addbibresource{literatur.bib}

\usepackage{listings}
\usepackage{xcolor}

\definecolor{codegreen}{rgb}{0,0.6,0}
\definecolor{codegray}{rgb}{0.5,0.5,0.5}
\definecolor{codepurple}{rgb}{0.58,0,0.82}
\definecolor{backcolour}{rgb}{0.95,0.95,0.92}

\lstdefinestyle{mystyle}{
    backgroundcolor=\color{backcolour},   
    commentstyle=\color{codegreen},
    keywordstyle=\color{magenta},
    numberstyle=\tiny\color{codegray},
    stringstyle=\color{codepurple},
    basicstyle=\ttfamily\footnotesize,
    breakatwhitespace=false,         
    breaklines=true,                 
    captionpos=b,                    
    keepspaces=true,                 
    numbers=left,                    
    numbersep=5pt,                  
    showspaces=false,                
    showstringspaces=false,
    showtabs=false,                  
    tabsize=2
}

\lstset{style=mystyle}

% Custom environments
\newenvironment{myitemize}{\begin{itemize}\itemsep -4pt}{\end{itemize}}
\newenvironment{myenumerate}{\begin{enumerate}\itemsep -4pt}{\end{enumerate}}
\newenvironment{myalphaenum}{\begin{enumerate}[label={\alph*)}]\itemsep -4pt}{\end{enumerate}}

% Custom commands
\newcommand{\absatz}{{\vspace{5mm}\newline}}
\newcommand{\zitat}[2]{
	\begin{quote}
		\textit{#1} \cite{#2}
	\end{quote}
}	
	
\newcommand{\geintrag}[2]{
	\textbf{#1}\begin{adjustwidth}{.5cm}{0cm}#2
	\end{adjustwidth}\vspace{3.5mm}
}
\newcommand{\singleenum}[2]{
	\begin{myenumerate}
		\setcounter{enumi}{#1}
		\item #2
	\end{myenumerate}
}

\newcommand{\specialcell}[2][l]{%
  \begin{tabular}[#1]{@{}l@{}}#2\end{tabular}}

% Settings for bibliography from BibLaTeX
\setcounter{biburllcpenalty}{7000}
\setcounter{biburlucpenalty}{8000}

\AfterTOCHead{\thispagestyle{empty}}

\begin{document}
\title{\Large{Untersuchung, Implementierungen und Bewertung von Graph-Metriken– und Ähnlichkeiten
}}
\author{Benedict Martin Weichselbaum}
\date{\today}
\selectlanguage{ngerman}
\begin{titlepage}
	\centering\hspace{8mm}
	\begin{figure}
		\centering
			\includegraphics[scale=1.3]{./Abbildungen/dhbwlogo.png}
	\end{figure}
	
	\vspace{5mm}	
	{\fontsize{26}{40}\selectfont
	Untersuchung, Implementierungen und Bewertung von Graph-Metriken– und Ähnlichkeiten
	}
	\\
	\vspace{2cm}
	\textbf{\Large{Studienarbeit}} \par
	\vspace{1cm}
	im Studiengang Informatik \par
	\vspace{0.3cm}
	an der Dualen Hochschule Baden-Württemberg Stuttgart, Campus Horb am Neckar \par
	\vspace{1.2cm}
	von \par
	\vspace{0.5cm}
	\textbf{\large{Benedict Weichselbaum}} \par
	\vspace{1.5cm}
	{\today}\par
	\vfill
	\begin{table}[ht]
		\hspace{1,5cm}
		\begin{tabular}{p{7cm}p{7cm}}
			\textbf{Bearbeitungszeitraum} & 28.09.2020 - 31.05.2021\\
			\textbf{Matrikelnummer, Kurs} & 6275457, TINF2018\\
			\textbf{Betreuer \& Gutachter} & Prof. Dr. ing. Olaf Herden\\
		\end{tabular}
	\end{table}
\end{titlepage}

\section*{Erklärung}
\thispagestyle{empty}
Ich versichere hiermit, dass ich meine Studienarbeit mit dem Thema \textit{Graphen: Metriken und Ähnlichkeit} selbstständig verfasst und keine anderen als die angegebenen Quellen und Hilfsmittel benutzt habe.
\newline
Ich versichere zudem, dass die eingereichte elektronische Fassung mit der gedruckten Fassung übereinstimmt.\vspace{1.6cm}\newline
{Nürnberg, \today\vspace{1.2cm}\par\vspace{1.5cm}}
{\noindent\rule{6cm}{.4pt}\newline Benedict Martin Weichselbaum}

\selectlanguage{english}
\begin{abstract}
	\cite{wills_metrics_2019}
\end{abstract}

\selectlanguage{ngerman}
\setcounter{tocdepth}{3}

{\tableofcontents \thispagestyle{empty}}


\listoffigures\thispagestyle{plain}
\pagenumbering{Roman}
\setcounter{page}{1}

\listoftables\thispagestyle{plain}

\chapter*{Abkürzungsverzeichnis}
\addcontentsline{toc}{chapter}{Abkürzungsverzeichnis}
\begin{acronym}

\end{acronym}


\chapter{Einleitung}
\pagenumbering{arabic}
\setcounter{page}{1}
\section{Motivation für die Studienarbeit}

\section{Fragestellungen}

\chapter{Graphmetriken und die Ähnlichkeit von Graphen}

\section{Verfügbare Graphmetriken}

\section{Verbreitung der vorgestellten Metriken}

\section{Bewertung der Metriken}

\section{Der Ähnlichkeitsbegriff bei Graphen}

\chapter{Untersuchung und Implementierung von Graphmetriken}

\section{Implementierung in verschiedenen Graphdatenbanken}

\section{Vergleich der Implementierungen}

\chapter{Graphmetriken und Ähnlichkeit in Anwendung}

\chapter{Fazit und Zusammenfassung}

\section{Zusammenfassung der Ergebnisse}

\section{Fazit}

\chapter*{Glossar}
\addcontentsline{toc}{chapter}{Glossar}
{
}
\printbibliography
\end{document}